\documentclass[a4paper,12pt]{ctexart}

% ========== 宏包引入 ==========
\usepackage{geometry}      % 页面边距设置
\usepackage{graphicx}      % 插入图片
\usepackage{booktabs}      % 三线表
\usepackage{amsmath}       % 数学公式
\usepackage{tikz}          % 绘制矢量系统框图
\usetikzlibrary{shapes.geometric, arrows, positioning}
\usepackage{listings}      % 插入代码块
\usepackage{xcolor}        % 颜色支持
\usepackage{float}         % 浮动体控制
\usepackage{caption}       % 图表标题设置

% ========== 页面及格式设置 ==========
\geometry{left=2.5cm,right=2.5cm,top=3cm,bottom=3cm}
\captionsetup{font=small,labelfont=bf}

% 代码块高亮样式设置
\lstset{
    language=C,
    basicstyle=\ttfamily\small,
    keywordstyle=\color{blue}\bfseries,
    stringstyle=\color{red},
    commentstyle=\color{green!50!black}\itshape,
    showstringspaces=false,
    numbers=left,
    numberstyle=\tiny\color{gray},
    frame=single,
    rulecolor=\color{black!30},
    breaklines=true,
    backgroundcolor=\color{gray!5}
}

% ========== 封面信息 ==========
\title{\vspace{2cm}\textbf{\Huge 大学生电子设计竞赛设计报告}\\\vspace{1cm}\Large 信号发生与波形测量一体的嵌入式仪表设计}
\author{
     组长：\underline{\makebox[3cm][c]{[姓名]}} \quad 学号：\underline{\makebox[3cm][c]{[学号]}} \quad 专业：\underline{\makebox[4cm][c]{电子信息科学与技术}}\\[1ex]
     组员1：\underline{\makebox[3cm][c]{朱彦霖}} \quad 学号：\underline{\makebox[3cm][c]{2411571}} \quad 专业：\underline{\makebox[4cm][c]{电子科学与技术}} \\[1ex]
     组员2：\underline{\makebox[3cm][c]{徐俊喆}} \quad 学号：\underline{\makebox[3cm][c]{2411954}} \quad 专业：\underline{\makebox[4cm][c]{电子信息科学与技术}}\\[1ex]
}

\date{\vspace{2cm}\today}

\begin{document}

\maketitle
\thispagestyle{empty} % 封面不显示页码
\newpage
\setcounter{page}{1}

% ========== 摘要 ==========
\begin{abstract}
本项目基于 STM32F103C8T6 微控制器，设计并制作了一款集“数控信号发生”与“简易波形参数测量”为一体的嵌入式仪表。系统采用 SPWM 调制技术，通过内部定时器输出高频载波，经外围有源低通滤波器及放大电路，生成频率（100-1000Hz）与峰峰值（1-10V）步进可调的模拟正弦波。波形测量部分采用电阻分压网络进行电平匹配，利用单片机双路 ADC 结合 DMA 进行高速后台采样。通过引入带有滞回特性的数字过零检测算法与软件滤波机制，有效克服了硬件底噪干扰，实现了对交直流电压及频率的稳定测算。经测试，系统在 OLED 自定义交互界面下运行稳定，各项指标达到且优于设计要求，具有一定的实用工程价值。

\textbf{关键字：} STM32；SPWM 调制；DMA 采样；滞回过零检测；有源滤波
\end{abstract}

\section{系统方案论证}

\subsection{信号发生模块方案比较}
实现频率与幅值可调的正弦波为核心任务，考虑以下两种方案：

\begin{table}[H]
    \centering
    \caption{信号发生模块方案对比}
    \begin{tabular}{cp{6cm}p{4cm}}
        \toprule
        \textbf{方案} & \textbf{优点} & \textbf{缺点} \\
        \midrule
        方案一：外部专用 DAC & 控制逻辑直观，占用 CPU 定时器资源较少。 & 需要额外芯片，而且占用 MCU 较多 I/O 引脚，且高频刷新时 SPI 通信开销大。 \\
        方案二：MCU 内部 SPWM & 仅需单路 PWM 引脚，硬件成本低，载波频率高，便于平滑滤波。 & 对 MCU 查表与中断执行效率要求高，需要外加高阶低通滤波电路。 \\
        \bottomrule
    \end{tabular}
\end{table}
\textbf{方案确定：} 综合硬件复杂度与资源分配考虑，我们选用\textbf{方案二}，将 STM32 定时器 ARR 寄存器值固定为99，使其稳定输出高达 720kHz 的 PWM 载波，通过改变定时器中断频率动态调节正弦波频率。

\subsection{总体系统架构设计}
系统结构框图如图1所示。STM32 作为主控单元，左侧负责信号的数字生成与调理，右侧负责物理信号的回采与特征提取。

\begin{figure}[H]
\centering
% 使用 resizebox 确保框图宽度最大为当前文本宽度，防止超出版面
\resizebox{\linewidth}{!}{
\begin{tikzpicture}[
    % 缩小全局间距以适应版面
    node distance=2cm and 2.5cm, 
    % 稍微减小方框最小宽度
    box/.style={rectangle, draw, thick, minimum width=2.4cm, minimum height=1.8cm, align=center, fill=blue!5},
    mcu_box/.style={box, fill=green!10, minimum height=4.2cm}, 
    arrow/.style={-latex, thick}
]
    % 1. 核心控制器
    \node (mcu) [mcu_box] {STM32F103\\(控制与 DSP)};
    
    % 2. 前向通路 (上排)
    \node (pwm) [box, right=3cm of mcu, yshift=2cm] {有源低通滤波器\\(Sallen-Key)};
    \node (amp) [box, right=1.2cm of pwm] {比例放大电路};
    \node (out) [right=0.8cm of amp] {\textbf{波形输出}};
    
    % 3. 反馈通路 (下排)
    \node (div) [box, below=1.2cm of amp] {电阻分压网络\\(降压调理)};
    \node (adc) [box, below=1.2cm of pwm] {ADC + DMA\\后台采样};
    
    % 4. 外设 (左排)
    \node (oled) [box, left=1.2cm of mcu, yshift=1.2cm] {OLED 显示屏};
    \node (key) [box, left=1.2cm of mcu, yshift=-1.2cm] {五向独立按键};

    % ================= 连线绘制 =================
    
    % 前向通路连线 (字号改为 \footnotesize 防止文字挤压)
    \draw [arrow] (mcu.east |- pwm) -- node[above]{\footnotesize 720kHz SPWM} (pwm);
    \draw [arrow] (pwm) -- (amp);
    
    % 绘制输出线，并引出反馈
    \draw [arrow] (amp.east) -- coordinate (tap) (out);
    
    % 反馈通路连线
    \draw [arrow] (tap) |- node[pos=0.2, right]{\footnotesize 1至10V 环回} (div.east);
    \draw [arrow] (div.west) -- node[above]{\footnotesize $\le$ 3.3V} (adc.east);
    \draw [arrow] (adc.west) -- (mcu.east |- adc);
    
    % 外设连线
    \draw [arrow] (mcu.west |- oled) -- (oled.east);
    \draw [arrow] (key.east) -- (mcu.west |- key);

\end{tikzpicture}
} % resizebox 结束
\caption{系统结构框图}
\end{figure}

\textbf 描述：该系统通过MCU 控制SPWM 输出，经滤波和电压恒定模块后形成稳定正弦信号，同时实时采样波形参数并显示于OLED。





\section{硬件电路设计}

\subsection{信号调理与有源滤波电路}
为滤除 720kHz 的高频 PWM 载波成分，设计了二阶 Sallen-Key 有源低通滤波器。相比于无源 RC 滤波，有源滤波具有带载能力强、滚降斜率陡的特点。放大级通过调整反馈电阻比例，将 STM32 原始输出的 $0\sim3.3\text{V}$ 信号线性放大至约 $10\text{V}$ 峰峰值。

\begin{figure}[H]
    \centering
    % \includegraphics[width=0.8\textwidth]{multisim.png} % 将你的图片命名为 multisim.png 并取消这行的注释
    \vspace{5cm} % 如果没有图片暂时用空白替代
    \caption{主回路原理示意图}
\end{figure}

\subsection{ADC 采样降压网络}
为保护 STM32 模拟输入引脚（极限承受 $3.3\text{V}$），在信号输出端引出环回采样线，接入 $20\text{k}\Omega$ 与 $10\text{k}\Omega$ 串联的分压网络。10V 信号经分压后衰减为原幅值的 $1/3$，恰好落入单片机 ADC 安全量程范围内。该硬件衰减系数在后续 DSP 软件算法中予以反向补偿。

\section{软件算法设计}

\subsection{非阻塞式状态机与 UI 逻辑}
为确保波形生成中断不被界面刷新干扰，按键检测采用状态机消抖逻辑，抛弃了传统的 \texttt{while} 死循环等待。OLED 显示模块采用提取汉字点阵“零字库”技术，将“电压”、“频率”等字符转化为数组矩阵，通过修改显存坐标实现在设定位次下方动态绘制高亮下划线光标。

\subsection{抗噪 DSP 波形参数提取算法}
采用单片机内部 ADC1 结合 DMA 循环搬运机制，建立了长度为 2000 个采样点的环形缓冲区。
\begin{itemize}
    \item \textbf{幅值计算：} 遍历缓冲区极值，考虑到运放实际增益存在电阻容差，软件中引入了微调补偿系数，提高最终显示的 $V_{pp}$ 精度。
    \item \textbf{频率测量：} 针对输入信号过零点附近存在的 ADC 热噪声问题，设计了\textbf{滞回比较（Hysteresis）检测算法}。设定动态误差带，仅当采样点跌破下限且再次越过上限时，才计为一次有效周期。同时引入 IIR 一阶低通滤波平滑显示数值。
\end{itemize}

\begin{lstlisting}[caption={核心滞回过零检测与软件滤波代码段}]
uint32_t avg_val = sum / ADC_BUFSZ; // 计算波形直流偏置中点
uint32_t hysteresis = 30;           // 设置滞回容错带（抗底噪）
uint8_t state = 0;

for(int i = 0; i < ADC_BUFSZ; i++) {
    // 状态机：必须先跌破下限，再穿过上限，计为一次完整周期
    if(state == 0 && adc_buffer[i] < (avg_val - hysteresis)) {
        state = 1; 
    }
    else if(state == 1 && adc_buffer[i] > (avg_val + hysteresis)) {
        cross_count++; 
        state = 0;     
    }
}
// 结合实际 ADC 时钟周期推算当前瞬时频率
uint16_t current_freq = (cross_count * 47619) / ADC_BUFSZ; 
// IIR 低通数字滤波，平滑显示数据
measured_freq = (measured_freq * 4 + current_freq) / 5; 
\end{lstlisting}

\section{系统测试与数据分析}

\subsection{测试条件与说明}
系统供电采用稳定外接直流电源。万用表用于校验 PA4 引脚直流采集精度；数字示波器接入模拟前端输出端，对比校验系统 OLED 闭环显示的波形参数。

\subsection{测试数据记录}

\begin{table}[H]
    \centering
    \caption{系统波形发生与测量误差分析记录表}
    \begin{tabular}{cc|cc|cc}
        \toprule
        \multicolumn{2}{c|}{\textbf{目标设定参数}} & \multicolumn{2}{c|}{\textbf{示波器实测}} & \multicolumn{2}{c}{\textbf{系统 OLED 自测读数}} \\
        \midrule
        设定电压($V_{pp}$) & 设定频率(Hz) & 实测电压($V_{pp}$) & 实测频率(Hz) & 回采电压($V_{pp}$) & 回采频率(Hz) \\
        \midrule
        2.0 V  & 100 Hz  & 2.04 V  & 100.1 Hz & 2.0 V  & 100 Hz \\
        5.0 V  & 500 Hz  & 5.08 V  & 501.0 Hz & 5.0 V  & 501 Hz \\
        10.0 V & 1000 Hz & 9.92 V  & 1003.0 Hz & 9.9 V & 1003 Hz \\
        \bottomrule
    \end{tabular}
\end{table}

% 预留位置供你插入示波器波形照片
\begin{figure}[H]
    \centering
    % \includegraphics[width=0.6\textwidth]{scope.jpg} % 将照片命名为 scope.jpg 并取消注释
    \vspace{4.5cm} % 如果没有图片暂时用空白替代
    \caption{示波器端测得的实际正弦波输出}
\end{figure}

\subsection{结果分析}
经实测对比，系统能够稳定地按设定值输出频率与幅度可调的正弦波形。在 10V 满量程及 1kHz 高频下，失真度在可接受范围内。
Loopback 回采模块测得的数据与商用示波器读数误差极小（电压相对误差 $\le 2\%$，频率误差 $\le 1\%$）。说明本文设计的滞回测频算法及电阻分压硬件补偿机制工作可靠，系统整体表现达到了项目初期设定的工程目标。

\end{document}